\documentclass[11pt]{article}
\usepackage[T1]{fontenc}
\usepackage{lmodern}
\usepackage[margin=0.9in]{geometry}
\usepackage{amsmath,amssymb,amsthm,mathtools}
\usepackage{booktabs,longtable,array,enumitem,microtype}
\usepackage{xcolor,fancyhdr}
\usepackage{xurl}
\usepackage[breaklinks,colorlinks=true,linkcolor=blue!45!black,citecolor=blue!45!black,urlcolor=blue!50!black]{hyperref}
\hypersetup{pdftitle={Smaller Diophantine consistency equations},pdfauthor={Prepared for Vladimir Reshetnikov},pdfsubject={Pell saturation, transport projection, and coprime proof-flow certificates},pdfkeywords={Diophantine equations, ZFC, consistency, Pell equations}}
\usepackage{listings}
\definecolor{ink}{RGB}{28,46,62}
\definecolor{accent}{RGB}{21,100,116}
\setlength{\headheight}{15pt}
\setlength{\parskip}{4pt}
\setlength{\parindent}{0pt}
\setlist{nosep,leftmargin=1.5em}
\pagestyle{fancy}
\fancyhf{}
\fancyhead[L]{\small\color{ink}Diophantine consistency: saturation and projection}
\fancyhead[R]{\small\texttt{c1a102936aca}}
\fancyfoot[C]{\thepage}
\lstset{basicstyle=\ttfamily\small,breaklines=true,columns=fullflexible,keepspaces=true,frame=single,rulecolor=\color{black!20}}
\newtheorem{theorem}{Theorem}[section]
\newtheorem{lemma}[theorem]{Lemma}
\newtheorem{proposition}[theorem]{Proposition}
\newtheorem{corollary}[theorem]{Corollary}
\theoremstyle{definition}
\newtheorem{definition}[theorem]{Definition}
\theoremstyle{remark}
\newtheorem{remark}[theorem]{Remark}
\newcommand{\Z}{\mathbb Z}
\newcommand{\N}{\mathbb N}
\newcommand{\Q}{\mathbb Q}
\newcommand{\Zp}{\mathbb Z_{>0}}
\newcommand{\Npos}{\mathbb Z_{>0}}
\newcommand{\Con}{\operatorname{Con}}
\newcommand{\ZFC}{\mathsf{ZFC}}
\newcommand{\Pos}{\operatorname{Pos}}
\newcommand{\file}[1]{\nolinkurl{#1}}
\newcommand{\pin}{\texttt{c1a102936aca437e9999530a0111e6ddab56dc6d}}
\allowdisplaybreaks[2]
\newcommand{\CheapDegree}{1,628,768}
\newcommand{\CheapOps}{434}
\newcommand{\CheapNaturalOps}{441}
\newcommand{\RadicalDegree}{1,103,950}
\newcommand{\RadicalOps}{1,074}
\newcommand{\RadicalNaturalOps}{1,081}


\begin{document}
\thispagestyle{empty}
{\sffamily\bfseries\color{accent}RESEARCH CONTRIBUTION\hfill 23 OCTOBER 2026\par}
\vspace{0.65cm}
{\Huge\bfseries\color{ink}Smaller Diophantine\\[4pt]consistency equations\par}
\vspace{0.35cm}
{\Large Pell saturation, exact transport elimination,\\and coprime proof-flow certificates\par}
\vspace{0.4cm}
Prepared for Vladimir Reshetnikov and the ProveIt project.\\
AI-assisted mathematical research with explicit polynomial artifacts.

\begin{abstract}
We improve two different parts of the current ProveIt frontier. A transport
coordinate change, already available in an earlier source but not applied to
the newest joint-sign family, gives one polynomial in \textbf{18 unrestricted
integers}, of exact degree \textbf{88}, evaluated in \textbf{105} binary
arithmetic operations. Its predecessor has 19 variables and 109 operations.
Three companion projections improve the variable counts at degrees 32, 20,
and 16 as well.

A separate Pell-certificate reduction removes a redundant factor from its
strong ordinate. The square root needed to restore that factor is already
certified by another part of the same equation. This gives a
\textbf{nine-positive-variable} polynomial of degree \textbf{\CheapDegree}
in \textbf{\CheapOps} operations, improving the current 440-operation,
degree-1,845,940 construction in both measures. A complete-radical majorant
gives another nine-variable point of degree \textbf{\RadicalDegree} in
\textbf{\RadicalOps} operations. All fixed coefficients and ordinary
division-free evaluation circuits are supplied.

The new existence equivalences are proved with explicit reverse maps. Their
consistency interpretation composes with the precisely pinned ordinary
native/compiler and syntactic-refutation calibration theorems. We do not
claim a new end-to-end Lean/Rocq certificate, a completed derivation in the
repository's encoded object-language ZFC, or global minimality. A distinct
Fermat-label proof-flow construction removes acyclicity and common-multiple
parameters at the finite-certificate level; its final fixed-dimensional
Diophantine closure remains open and is not counted as a finished equation.
\end{abstract}
\vfill
\small
Reviewed repository snapshot: \pin.\\
Both supplied repository directories were inspected, including all five
archives then in \file{docs/incoming}. All comparisons are to this frozen
snapshot, not to a moving branch. The local date is 23 October; the snapshot's
commit timestamp is 24 October in UTC.
\newpage
\normalsize
\tableofcontents
\newpage

\section{The mathematical target and the actual improvement}

\subsection{A single fixed polynomial and a specified domain}

Fix the recursive proof predicate of the syntactic ZFC calculus used in
ProveIt. With $\bot$ its chosen contradictory sentence, write
\[
 \Con(\ZFC)\;:=\;\forall p\in\N\;
          \neg\operatorname{Proof}_{\ZFC}(p,\ulcorner\bot\urcorner).
\]
The requested endpoint is an integer-coefficient polynomial $P$, with no
remaining program or theory parameter, and a domain $D$ for which
\begin{equation}\label{eq:target}
 \ZFC\vdash
 \bigl(\neg\exists\mathbf x\in D^m\ P(\mathbf x)=0\bigr)
       \longleftrightarrow\Con(\ZFC).
\end{equation}
The two sides of this biconditional are not asserted separately. A positive
solution, a natural solution, and an unrestricted integer solution are
different assertions. In this report $\N$ includes zero.

All displayed intermediate symbols are polynomial abbreviations unless
expressly declared as unknowns. Fixed powers have fixed nonnegative integer
exponents. Fixed coefficients are given both as lossless integer bytes and
by deterministic finite recipes. No coefficient depends on a consistency
decision, a future search, or a choice of a refutation. The literal JSON
circuits define ordinary polynomials even where their full monomial
expansions would be inconvenient to print.

\subsection{Measured results}

\begin{center}
\small
\begin{tabular}{@{}llrrr@{}}
\toprule
Construction & Domain & Variables & Exact degree & Operations\\
\midrule
Transport projection & $\Z$ & 18 & 88 & 105\\
Companion & $\Z$ & 19 & 32 & 108\\
Companion & $\Z$ & 20 & 20 & 111\\
Companion & $\Z$ & 21 & 16 & 114\\
Pell saturation, short bound & $\Zp$ & 9 & \CheapDegree & \CheapOps\\
Pell saturation, radical bound & $\Zp$ & 9 & \RadicalDegree & \RadicalOps\\
\bottomrule
\end{tabular}
\end{center}

For each of the first four rows, its predecessor has one more integer
variable, the same degree, and four more operations. The reduction is an
exact integer-polynomial projection, so it applies to every source zero;
it does not rely on a new restriction to specially chosen canonical zeros.

The short nine-positive construction reduces degree from 1,845,940 to
\CheapDegree, an improvement of approximately 11.76\%, while also reducing
440 operations to \CheapOps. Its natural-number translation has the same
degree and \CheapNaturalOps\ operations. The complete-radical construction
is a separate degree/circuit tradeoff; the inspected incoming radical
candidate has degree 1,158,228 and 1,753 operations. The new point has lower
degree and the measured circuit cost given above. The smaller ten-integer
and nine-natural variable counts already present in the repository remain
separate results: the 18-integer equation is useful for its much lower degree
and short expression, not as a fewest-variable record.

An operation is a live binary addition, subtraction, or multiplication.
Squaring and multiplication by a fixed scalar are charged. Reusing an
already computed value and supplying a fixed integer leaf are free under
the repository's convention. This is arithmetic circuit size, not bit
complexity; the fixed coefficients are large. A sum of six residual squares
is one polynomial equation, rather than six separate asserted equations.

\subsection{What is new, and what is reused}

The transport coordinate change itself is already explicit in source 73,
\emph{Reanchored sign constraints and adaptive Pell carriers}. The new
result is its application to the newer joint-sign projection family, with
complete equations, inverse maps, exact degrees, and paid circuits. The
argument is elementary; the improved endpoint is the contribution.

The saturated positive Pell certificate changes the certificate itself.
It exploits a square root already forced elsewhere, retains the necessary
coprimality conditions, and replaces the old majorant. The incoming
round-eleven Lucas construction uses a related saturation idea in a
different unrestricted-integer certificate. We credit that precursor.
The positive certificate, its interface proof, and its full-equation
costs are established here. The subsequent complete-radical optimization
also builds on the incoming radical-majorant methodology; it is supported
by exact integer coefficient certificates, not an asymptotic expansion
used as an inequality.

The Fermat-label appendix is a genuinely different source architecture
relative to the inspected budget and height-register constructions. Its
ingredients are classical. We claim its proved finite-certificate theorem
and identify the remaining compiler costs; we do not claim historical
priority or an unmaterialized final polynomial record.

\subsection{Proof layers and formalization status}\label{sec:boundary}

There are four distinct kinds of evidence in this package:
\begin{enumerate}
 \item Complete ordinary proofs of the new projections, Pell saturation,
       radical bounds, and finite proof-flow equivalence.
 \item Exact finite checks binding the literal equations to those proofs:
       polynomial identities, coefficient reconstruction, nonzero leading
       coefficients, topology, liveness, and arithmetic operation counts.
 \item The inherited ordinary compiler/refutation calibration, pinned to
       immutable source data and retained proof notes. This supplies the
       link from the unchanged arithmetic source to ZFC refutations.
 \item A possible future kernel-checked derivation of the entire composed
       sentence inside the repository's encoded ZFC calculus.
\end{enumerate}
The new results establish the stated solvability reductions. Their
$\Con(\ZFC)$ interpretation follows by composition with the retained ordinary
calibration theorems. This contribution does not independently reprove the
entire inherited semantic chain or supply its unfinished kernel-checked
internalization. The repository itself records unfinished native
phase/global calibration and internal-model transport obligations; this
package does not silently mark them complete. Its already checked generic
consistency equation remains a separate formal checkpoint.

The elementary constructions use finite arithmetic, induction, Pell
recurrences, and finite coded derivations. They are suitable for
formalization in ZFC. To interpret \eqref{eq:target} as an ordinary theorem,
compose them with the named ordinary calibration proofs. To interpret it
as a claim about a newly exported proof object in the repository's formal
calculus would be stronger and is not justified by this delivery. Finite
tests alone do not prove an unbounded consistency equivalence.

\section{Removing one integer witness from the joint-sign family}
\label{int-section}

The smallest-degree branch and the smallest-variable branch need different
comparisons.  Here we improve four recent equations over \emph{unrestricted
integers}.  The principal endpoint is a single polynomial in eighteen
integer variables, of exact degree $88$, with $105$ live binary arithmetic
operations.  Its direct predecessor in the incoming joint-guards packet
uses nineteen variables, degree $88$, and $109$ operations.

\begin{center}
\begin{tabular}{rrrrrr}
\toprule
\multicolumn{3}{c}{Incoming predecessor}&\multicolumn{3}{c}{This composition}\\
Variables&Degree&Operations&Variables&Degree&Operations\\
\midrule
19&88&109&18&88&105\\
20&32&112&19&32&108\\
21&20&115&20&20&111\\
22&16&118&21&16&114\\
\bottomrule
\end{tabular}
\end{center}
Every entry in this table has domain $\Z$; none is obtained by silently
relabelling positive witnesses.  All four new equations retain precisely
the same six fixed numerals as their direct predecessors.  These are new
compositions and measured endpoints, not a new general elimination lemma:
source~73, \emph{Guard reanchoring}, already uses the complement coordinate
and monic transport elimination proved below.  Its application there is to
an earlier forty-one-variable system.  The incoming joint-sign projection
family retains the eliminable transport coordinate, which lets the same
method produce the four improvements in the table.

\subsection{An exact integer-polynomial projection}
\label{int-projection}

Suppose a source system contains polynomial abbreviations $q,W$ independent
of two integer coordinates $F,\alpha$, and put $C=q-F-\alpha$.
Suppose one of its rows is
\begin{equation}
 (K_0+W)C-F-z(q-1)=0.                         \label{int-old-transport}
\end{equation}
First replace $\alpha$ by the free integer coordinate $C$.  This is an
invertible polynomial change of coordinates, with inverse
$\alpha=q-F-C$.  The row is now monic in $F$.  Define
\begin{equation}
 \widehat F=(K_0+W)C-z(q-1),\qquad
 \widehat\alpha=q-\widehat F-C.               \label{int-restore}
\end{equation}
Substitute these polynomials in every other row and delete
Equation~\eqref{int-old-transport}.

\begin{proposition}
\label{int-projection-thm}
The old and new integer zero sets are in bijection.  The projection fixes
every other coordinate and sends $(F,\alpha)$ to $C=q-F-\alpha$;
its inverse is Equation~\eqref{int-restore}.
\end{proposition}
\begin{proof}
At an old zero the transport row gives $F=\widehat F$, and then
$\alpha=\widehat\alpha$.  All substituted rows vanish.  Conversely,
restore these two values at a new zero.  The deleted row vanishes
identically, while every retained row vanishes by the defining
substitution.  Both compositions fix the original coordinates, because
the old transport equation determines $F$ uniquely.  There is no division,
positivity assumption, coprimality assumption, or new choice of canonical
Pell witnesses in this argument.
\end{proof}

The rank expression also becomes shorter:
\begin{align}
 &\bigl((q-1)(q-F)+\alpha\bigr)(q^2-1)
       +(M_C+qM_F)J \notag\\
 &\hspace{15mm}=
   \bigl(q(q-\widehat F)-C\bigr)(q^2-1)
       +(M_C+qM_F)J.                         \label{int-rank-identity}
\end{align}
This is a literal polynomial identity after the substitution, not merely
an equality at solutions.  The determinant guard remains valid with
$F=\widehat F$ because it is substituted everywhere, including its sign
constraints.  No coefficient port is replaced or borrowed from another
construction.

\subsection{The complete eighteen-variable polynomial}
\label{int-polynomial}

The eighteen unknowns are exactly
\begin{equation}
 J_3,C,z,f,h,r,v,w,\tau,y,\gamma,o,j,c,t,p_1,p_2,p_3\in\Z.
                                                        \label{int-vars}
\end{equation}
The variable $C$ is unrestricted; its required sign is enforced by the
polynomial's inherited guard.  All following symbols are polynomial
abbreviations and do not add quantified variables.  Define
\begin{align*}
 J&=3J_3+1,&q&=\beta J+1,\\
 W&=4(3w+1),&X&=qW,\\
 Y&=(2v+1)q^3,&F&=(K_0+W)C-z(q-1),\\
 E&=XY,&a&=E+Y,\\
 H&=4a+3,&\Delta&=a^2+H,\\
 V&=fo-c,&I&=(2r+1)c^2,\\
 k&=1-V+jc+hE,&\eta&=c-kY,\\
 \zeta&=k-\eta,&T&=\eta\zeta.
\end{align*}
Then put
\begin{align*}
 R&=\bigl(q(q-F)-C\bigr)(q^2-1)+(M_C+qM_F)J,\\
 d&=(A F)(B C)-(q-AF-BC)^2,\\
 P_+&=1+p_1^2+p_2^2+p_3(p_3+1).
\end{align*}
Here $P_+\ge1$ for every integer triple; its standard surjectivity onto
positive integers is used only in the inherited joint-sign theorem.
The six residuals are
\begin{align}
 L_1&=\tau^2-EkY(k+EkY)-1,\notag\\
 L_2&=(X+ac+\gamma H)^2-\Delta c^2-1,\notag\\
 L_3&=cj-V-R,\notag\\
 L_4&=I^2-\Delta(f^2-1),\notag\\
 L_5&=I^2(V^2-y^2)+y^2-1,\notag\\
 L_6&=dT-(t^2-d-T+1)^2-P_+.                    \label{int-six-rows}
\end{align}
The complete equation is
\begin{equation}
 \boxed{\quad Q_{18}=L_1^2+L_2^2+L_3^2+L_4^2+L_5^2+L_6^2=0,
          \qquad\mathbf x\in\Z^{18}.\quad}    \label{int-Q18}
\end{equation}
The JSON uses the names \texttt{v\_scale}, \texttt{w3},
\texttt{pair\_root}, and \texttt{pair\_pos1--3} for $v,w,t,p_1,p_2,p_3$.
Those are only spelling changes.  Its $105$-gate straight-line program is
an explicit division-free definition of the same ordinary polynomial.

\subsection{The six fixed numerals}
\label{int-coefficients}

The tuple is the complete source-34 native binding named
\texttt{mixed-parity-unrounded-1064921}, together with the two determinant
scales.  In the displayed equation,
\[
 \beta=2^{1064921}-1,\qquad A=2^{1113},\qquad B=2^{666104}.
\]
The other three fixed integers $K_0,M_C,M_F$ are supplied as literal,
compressed, minimal unsigned little-endian byte strings, with their
complete finite construction retained alongside them.  In particular
$M_F$ is the \emph{already shifted} source port.  It is not shifted again.
All six numeral payloads together occupy $615883$ bytes before compression.

For clarity, the finite construction of the native ports is as follows.
The retained literal $215$-state table determines $943$ ordered Wang tiles
by the explicit finite map in
\path{integer/upstream/guards/upstream/source34/source34_builder.py}.
With radix $\rho=2^{121}$ and block length $8801$, that map builds a
$989$-element support $\mathcal E$, its explicit nonnegative coefficients
$c_e$, and the dictionary
\[
 d_s=\mathbf1_{s\in\{40,41,1100,2200\}}
      +\sum_{e\in\mathcal E}c_e
          \bigl(\mathbf1_{s=4393-e}+\mathbf1_{s=5495-e}\bigr).
\]
Here the indicator sum counts coincident contributions.  Put
$D_C=\sum_s d_s\rho^s$ and $\mathcal B=\rho^{8801}$.
The finite recipes then give
\begin{align*}
 K_0&=D_C+\mathcal B\rho,\\
 M_C&=\mathcal B-1
         -\sum_{e\in\mathcal E\setminus\{1\}}\rho^e
         -2\rho^{1063},\\
 M_F&=(\mathcal B-1)+4+
      \mu(\rho^{4393}+\rho^{5495})
      +\sum_{j=0}^{9}(\rho^{1000+2j}+\rho^{1021+2j})\\
     &\hspace{8mm}+\rho^{2195}+\rho^{3296}
       +(2^{119}-1)\sum_{e=8792}^{8798}\rho^e
       +(2^{27}-1)\rho^{8799},\\
 \mu&=1022+\sum_{j=0}^{43}2^{10+2j}.
\end{align*}
The finite support, every tile and coefficient, and the generating table
are included as data and executable finite recipes.  This is not an
appeal to an unspecified universal polynomial or a free compiler index.
The verifier reconstructs the entire tuple, compares its integer values
to the literal bytes, and replays $69230$ finite convolution-basis checks.
These finite checks bind the numerals; the unbounded semantic theorem is
a separate inherited argument.

\subsection{Why the joint row really enforces two positive signs}

The inherited guard has a useful elementary proof, which also explains why
the eighteen displayed inputs may range over all integers.  For integers
$x,y,t$ put
\[
 G(x,y,t)=xy-(t^2-x-y+1)^2.
\]
Then $x,y>0$ if and only if $G(x,y,t)>0$ for some integer $t$.
Indeed, $G>0$ implies $xy>0$.  If $x=-a,y=-b$ with $a,b>0$, then
$(t^2+a+b+1)^2>ab$, a contradiction.  Conversely, for positive $x,y$ put
$n=x+y-1$.  We have $xy\ge n$.  If $n$ is a square, take $t^2=n$.
Otherwise write $m^2<n<(m+1)^2$ and choose whichever of $m,m+1$ makes
$|t^2-n|$ smaller.  This distance is at most $m$, so
$(t^2-n)^2\le m^2<n\le xy$.

Every positive integer is a value of
$1+p_1^2+p_2^2+p_3(p_3+1)$.  This follows from the three-square theorem
applied to $4n-3$: exactly one square is odd modulo four, giving the
displayed expression after division by four.  Consequently
\[
 x,y>0\quad\Longleftrightarrow\quad
 \exists t,p_1,p_2,p_3\in\Z\quad
 G(x,y,t)=1+p_1^2+p_2^2+p_3(p_3+1).
\]
Apply this to $x=d,y=T$ in $L_6$.  The new projection substitutes $F$
in both occurrences of the determinant and in the rank; it does not
delete this guard or infer positivity from the spelling of a variable.
The further interpretation of $d>0$ and $T>0$ is exactly the retained
parent's determinant and signed-Pell argument.

\subsection{Correctness and its precise proof boundary}
\label{int-correctness}

Let $Q_{19}^{\mathrm{old}}$ be the exact incoming polynomial
\path{joint/joint19_degree88.json} in the archive
\path{ZFC_Diophantine_Joint_Guards_2026-10-24.zip}.
The repository is pinned at commit \texttt{c1a102936aca};
the complete identifier appears on the title page.
Its source graph and the exact parent-to-polynomial binding are retained.
The preceding projection gives the unconditional identity
\[
 Q_{18}(\mathbf x)=
 Q_{19}^{\mathrm{old}}(\mathbf x,
       F=\widehat F,\alpha=\widehat\alpha),
\]
where the input correspondence replaces the old $F,\alpha$ by the new
$C$.  The deleted squared residual is identically zero.
Consequently
\begin{equation}
 \exists\mathbf x\in\Z^{18}\ Q_{18}=0
 \quad\Longleftrightarrow\quad
 \exists\mathbf y\in\Z^{19}\ Q_{19}^{\mathrm{old}}=0.       \label{int-bijection}
\end{equation}
This equivalence is proved over arbitrary integer values of the six
coefficient ports, so there is no hidden numerical or sign assumption
in the new step.

For the authentic tuple the retained parent proof proceeds through the
joint-positivity lemma, determinant sign restoration, the signed-power
bootstrap, native Wang-tableau compilation, Pell calibration, and the
literal machine/refutation theorem.  Composing that ordinary proof chain
with Equation~\eqref{int-bijection} yields the intended theorem
\[
 \ZFC\ \vdash\
 \bigl(\neg\exists\mathbf x\in\Z^{18}\ Q_{18}(\mathbf x)=0\bigr)
       \ \longleftrightarrow\ \Con(\ZFC).
\]
This statement neither assumes nor proves consistency.  The new
projection is an elementary argument formalizable in ZFC.  The archive
does not provide a new end-to-end Lean/Rocq certificate or an exported
object-language ZFC derivation: it inherits the named unbounded
calibration interfaces, their ordinary proofs, and their existing formal
verification limitations.  A finite replay cannot replace those interfaces.

\subsection{Exact degree and operation count}
\label{int-degree}

Give each quantified coordinate degree one and the fixed ports degree
zero.  The six residual degrees in Equation~\eqref{int-six-rows} are
\[
 (34,14,5,14,10,44).
\]
In particular, the rank row rises from degree four to degree five after
the projection; this does not change the dominant degree.  Write
$X_*,Y_*,E_*,\eta_*$ for the top homogeneous parts.  They are
\begin{align*}
 X_*&=36\beta J_3w,&Y_*&=54\beta^3J_3^3v,\\
 E_*&=X_*Y_*,&\eta_*&=-hX_*Y_*^2.
\end{align*}
Thus $T$ has degree $22$ and top part $-\eta_*^2$, whereas $d$ has
degree at most four.  Only $L_6^2$ reaches degree $88$, with top part
$\eta_*^8$.  In particular
\begin{equation}
 [h^8w^8v^{16}J_3^{56}]Q_{18}
       =2^{32}3^{64}\beta^{56}>0.                    \label{int-top}
\end{equation}
This proves exact degree on the actual fixed-coefficient slice, rather
than merely giving a degree propagated through a circuit.

The delivered DAG uses $54$ multiplications, $27$ additions, and $24$
subtractions.  Every square and every scalar multiplication is charged;
fixed numeral leaves and reuse of a register are free.  The six residuals
need $94$ operations, and six squarings plus five additions give $105$.
The predecessor needs $109$.  Restoring $F$ costs additional arithmetic,
but eliminating its transport row and shortening the rank expression
save four live operations overall.  No minimum-operation claim is made.

\subsection{Three companion tradeoffs}
\label{int-companions}

The same transformation applies to the incoming $20/32$, $21/20$, and
$22/16$ joint-sign polynomials.  Their three resulting packets have,
respectively, nineteen, twenty, and twenty-one unrestricted coordinates.
They retain some intermediate coordinates that Equation~\eqref{int-Q18}
eliminates.  Precisely:
\begin{itemize}
\item The nineteen-variable, degree-$32$ form retains $\eta$ and the row
      $c-\eta-kY=0$.
\item The twenty-variable, degree-$20$ form also retains $E$ and the row
      $E-XY=0$, while $a=E+Y$ remains an abbreviation.
\item The twenty-one-variable, degree-$16$ form also retains $a$ and
      $a-E-Y=0$.  Its source rows use $a-E$ in place of $Y$ where specified
      by the parent.  It retains the parent's reduced auxiliary residual
      $\Delta(f^2-1)(V^2-y^2)+y^2-1$.
\end{itemize}
The last replacement is justified in the parent by retaining
$I^2-\Delta(f^2-1)$ and subtracting its multiple $V^2-y^2$ from the
original auxiliary row.  It is a zero-set-preserving row operation, not
an unrestricted polynomial identity.  Our new substitution preserves the
particular parent polynomial for each endpoint.

Their independently checked degree certificates are
\begin{align*}
 [\eta^4h^4w^4v^4J_3^{16}]Q_{19}&=2^{12}3^{20}\beta^{16},\\
 [f^4v^4J_3^{12}]Q_{20}&=2^43^{12}\beta^{12},\\
 [a^4f^8o^4]Q_{21}&=1.
\end{align*}
All are positive for the authentic tuple.  The expanded row degrees and
literal circuit counts are in the verification receipts.  These exact
integer-domain improvements coexist with lower-variable, much-higher-degree
equations; no scalar measure makes every tradeoff comparable.

\subsection{Reproduction and provenance}
\label{int-reproduction}

The component commands are \texttt{python3 build.py},
\texttt{python3 verify.py}, and \texttt{python3 -O verify.py}.
The independent checker does not import the new builder.  It authenticates
the retained source files and coefficient tuple, rechecks the parent DAGs
against their literal polynomials, proves every new residual identity by
exact sparse algebra over validated cut registers, checks degree bounds
and nonzero top coefficients, verifies the sum-of-squares finalizer,
and rejects an operation-preserving sign corruption.  Normal and optimized
Python runs pass.  No diagnostic search for a consistency-equation zero
is used.

The primary polynomial is \path{transport18_degree88.json}; the three
companions are \path{transport19_degree32.json},
\path{transport20_degree20.json}, and
\path{transport21_degree16.json}.  Each declares every witness, fixed
numeral, gate, residual, and restoration map.  The inherited proof chain,
source~73's earlier method, literal finite recipes, fixed bytes, licenses,
and source hashes are retained under \texttt{upstream/}.  The method is
credited to that earlier report; the current contribution is its exact
composition with the newer joint-sign family.

\section{Removing a redundant Pell factor}
\label{nine-section}

The central improvement replaces a square of an already certified square
by that square itself.  The square root needed to explain the replacement
exists only in the proof; it does not become an additional unknown.
This changes the arithmetic certificate, so it requires a new soundness
argument.  It produces two complete division-free polynomial circuits:
\begin{center}
\begin{tabular}{lrrr}
Construction & Positive unknowns & Exact degree & Operations\\\hline
Frozen H2082 shared-product parent &9&1,845,940&440\\
Saturated certificate, simple bound &9&1,628,768&434\\
Saturated certificate, whole radical &9&1,103,950&1,074
\end{tabular}
\end{center}
Operations are live binary additions, subtractions and multiplications;
squares and scalar products are charged, and fixed numeral leaves and
register reuse are free.  The natural-number translations use respectively
441 and 1,081 operations, with unchanged degrees.  These are two different
tradeoffs, rather than a claim that either is globally simplest.

\subsection{The new Pell representation}

Let $\psi_A(0)=0$, $\psi_A(1)=1$ and
$\psi_A(t+2)=2A\psi_A(t+1)-\psi_A(t)$.  Let $\chi_A$ be the companion
sequence with initial values $1,A$.  Thus
\[
 \chi_A(t)^2-(A^2-1)\psi_A(t)^2=1.
\]
For $A\ge2$, the positive Pell solutions are exactly these pairs.

\begin{theorem}[Saturated positive Pell certificate]
\label{nine-saturation}
Let $A,B,C,J$ be positive integers satisfying
$A\ge3$, $1<B<A$, $C>B$ and $J\ge A$.  For $i\ge1$ and $j_0\ge0$ put
\begin{align*}
 D&=1+(A^2-1)C^2,&e_0&=iJC^2,\\
 F&=1+(A^2-1)De_0^2,&\Lambda&=CD+1-A,\\
 G&=A+F\Lambda,&H&=C+F(B+j_0C),\\
 I&=1+(G^2-1)H^2.&&
\end{align*}
Then $C=\psi_A(B)$ if and only if some such $i,j_0$ make $DFI$ a square.
\end{theorem}

\begin{proof}
All displayed radicands and $G,H$ are positive.  In particular
$D\ge A^2$ and $\Lambda\ge A^2-A+1>0$.  The congruences
\begin{align*}
 F&\equiv1\pmod D,&G&\equiv1\pmod D,&I&\equiv1\pmod D,\\
 G&\equiv A\pmod F,&H&\equiv C\pmod F,&I&\equiv D\pmod F
\end{align*}
show that $D,F,I$ are pairwise coprime.  If their product is a square,
each factor is a square.  Set $d=\sqrt D\in\Z_{>0}$.  Then
\[
 F=1+(A^2-1)(iJC^2d)^2.
\]
Pell classification gives positive indices $s,m,t$ such that
\[
 C=\psi_A(s),\qquad iJC^2d=\psi_A(m),\qquad H=\psi_G(t).
\]
The standard rank argument gives $C\mid m$.  Here are its essential
details, including the case of even $C$.  From
$\psi_A(s)\mid\psi_A(m)$ the divisibility law gives $s\mid m$; write
$m=su$.  Expanding the Pell addition identities yields
\[
 \psi_A(su)\equiv
 u\chi_A(s)^{u-1}\psi_A(s)\pmod{\psi_A(s)^3}.
\]
Since $\gcd(\chi_A(s),\psi_A(s))=1$, the stronger divisibility
$\psi_A(s)^2\mid\psi_A(m)$ implies $\psi_A(s)\mid u$, hence $C\mid m$.
The usual divisibility law itself follows by Euclidean reduction of the
addition identities and strict growth of $\psi_A$.

Because $G\equiv A\pmod F$ and $H\equiv C\pmod F$, recurrence reduction
gives
\[
 \psi_A(s)\equiv\psi_A(t)\pmod{\chi_A(m)}.
\]
This implies $s\equiv\pm t\pmod{2m}$.  A short proof avoids any
exceptional-index assumption.  Put $x=\chi_A(m)$.  The addition identities
give
\[
 \psi_A(v+2m)\equiv-\psi_A(v),\qquad
 \psi_A(2m-v)\equiv\psi_A(v)\pmod x.
\]
Every index folds to $r\in[0,m]$, with residue $\pm\psi_A(r)$ and index
congruent to $\pm r\pmod{2m}$.  The Pell identity and $A\ge3$ imply
$x>2\psi_A(m)$.  The signed representatives therefore lie strictly between
$-x/2$ and $x/2$.  Equal residues are equal integers, and strict growth
forces equal folded indices.

On the other hand $G\equiv1\pmod C$ and $H\equiv B\pmod C$, so
$t\equiv B\pmod C$.  Therefore $s\equiv\pm B\pmod C$.  Since $C>1$,
we have $s\ge2$; the inequality $\psi_A(s)\ge As$, together with $B<A$,
gives $0<s+B<C$.  Consequently the only possibility is $s=B$.

Conversely, suppose $C=\psi_A(B)$ and put $d=\chi_A(B)$, so $d^2=D$.
The determinant-one matrix
\[
 M_A=\begin{pmatrix}2A&-1\\1&0\end{pmatrix}
\]
has finite order modulo $JdC^2$.  Choose a positive multiple $m$ of this
order and set $i=\psi_A(m)/(JdC^2)>0$.  Then $F=\chi_A(m)^2$.
Let $T$ be a positive period of $M_G$ modulo $F$.  Choose arbitrarily large
$t\equiv B\pmod{TC}$ and put $H_0=\psi_G(t)$.  We have
\[
 H_0\equiv C\pmod F,\qquad H_0\equiv B\pmod C,
\]
using respectively $G\equiv A\pmod F$ and $G\equiv1\pmod C$.
Since $F\equiv1\pmod C$, the quotient $(H_0-C)/F$ is $B$ modulo $C$.
Taking $t$ sufficiently large makes that quotient at least $B$; hence
\[
 j_0=\frac{(H_0-C)/F-B}{C}\in\N.
\]
Then $H=H_0$ and $I=\chi_G(t)^2$.  This supplies the required witnesses.
All divisions here occur in the witness proof, not in the polynomial.
\end{proof}

The old certificate used
\[
 F_{\rm old}=1+(A^2-1)D^2(i_{\rm old}JC^2)^2.
\]
An old valid witness gives a particularly simple new witness:
$i=\sqrt D\,i_{\rm old}$.  It preserves $F,G,H,I$ exactly.  The converse
is supplied by Theorem~\ref{nine-saturation}, rather than by an unjustified
division of the free integer $i$ by $\sqrt D$.

\subsection{Binding to the actual compact54-sharp source}

Both new equations retain the entire H2082 source prefix and its coherent
compact54-sharp coefficient tuple from commit
\texttt{c1a102936aca437e9999530a0111e6ddab56dc6d}.
The complete layout, three kernel polynomials and coefficient leaves are
included as literal data.  The four original coefficient ports occupy
215,492 raw bytes; the eleven expanded live fixed leaves occupy 1,565,018
raw bytes.  These are different representations of the fixed data and
must not be confused with the number of unknowns.

For the inherited two-block words $S_{\rm code},T_{\rm code}$, define
\begin{align*}
 B_{\rm code}&=2^{864026}b^2,&
 L&=bB_{\rm code}^{2200},\\
 N&=2^{6480184}B_{\rm code}^{1162},&
 R&=L(S_{\rm code}+T_{\rm code}+1)-S_{\rm code}-1,\\
 t_{\rm cap}&=bB_{\rm code}^{116}.&&
\end{align*}
The supplied source tables define both words without a free program index.
The inherited source theorem gives $R,N\ge8$ and the required
carry/binomial calibration.  Define the ratio registers
\begin{align*}
 U&=N^2w,&Y&=N^2s,&M&=RY,\\
 P_0&=2M^2U,&K&=R+1+h(P_0-1),&A&=M(U+1),\\
 B_1&=2R+1,&W_1&=bw,&J&=4A-5,\\
 Q_0&=3W_1C-2(W_1^2-1),&&&\\
 a_0&=1+(P_0^2-1)K^2,&V_0&=P_0K,&&\\
 m_0&=K^2(4t_{\rm cap}-g)-16t_{\rm cap}(C-KY)^2,&&&\\
 \Omega&=Jm_0,&E_0&=J(n+W_1^2)+Q_0.&&
\end{align*}
The nine inputs are
\[
 b,g,h,s,w,C,i,j,n\in\Z_{>0}.
\]
Set $j_0=j-1$ and use $B=B_1$ in Theorem~\ref{nine-saturation}.
The source bounds imply $A\ge33280$, $1<B_1<A$ and $J\ge A$ on every
ambient positive assignment.  The hypothesis $C>B_1$ will be established
from the final equation before the representation theorem is invoked.

\subsection{The 434-operation equation}

Define
\[
 T_0=ADe_0,\qquad V_1=T_0GH,\qquad a_1=DFI,\qquad
 \Theta=E_0-\Omega(V_0+V_1).
\]
The first new equation is
\begin{equation}
 \boxed{\bigl(\Theta^2-\Omega^2(a_0+a_1)\bigr)^2
                    -4\Omega^4a_0a_1=0.}
 \label{nine-cheap-equation}
\end{equation}
Its unconditional majorant is exceptionally small:
\begin{align}
 T_0^2-DF&=D(De_0^2-1)>0,\label{nine-cheap-defect}\\
 V_1^2-a_1&=(T_0^2-DF)(GH)^2+DF(H^2-1)>0.\nonumber
\end{align}
Thus $V_1>\sqrt{a_1}>0$.  Also $V_0^2-a_0=K^2-1>0$.
These inequalities hold before a solution has been assumed.

\begin{proposition}
\label{nine-cheap-equivalence}
Equation~\eqref{nine-cheap-equation} has a positive solution exactly when
the same pinned native source has a solution.
\end{proposition}

\begin{proof}
The identity
\[
 E_0=Jn+(J-2)W_1^2+3W_1C+2>0
\]
shows that a zero cannot have $m_0=0$, since then its left side would be
$E_0^4$.  Factoring the two-root norm over the reals gives signed roots
$r_0^2=a_0$, $r_1^2=a_1$ satisfying
\begin{equation}
 \frac{E_0}{J}=m_0(V_0+V_1+r_0+r_1).
 \label{nine-sign-first}
\end{equation}
The strict majorants make the parenthesis positive for every sign choice;
therefore $m_0>0$ before any square-root integrality is used.

The unchanged margin argument now gives $|C-KY|<K/2$ and
$C>B_1+W_1$.  Moreover $A^2>P_0$, $D>C^2$, $F\ge A^2$ and $I\ge A^2$,
so
\[
 DFI>C^2A^4>K^2P_0^2>a_0.
\]
Hence the rational sum $x=\Theta/\Omega=r_0+r_1$ is nonzero.  The identity
$2xr_0=x^2+a_0-a_1$ makes each root rational, hence integral.
Equation~\eqref{nine-sign-first} then gives $J\mid Q_0$.
Theorem~\ref{nine-saturation} gives $C=\psi_A(B_1)$.
The retained arithmetic ratio/binomial theorem recovers the source zero.
If its original auxiliary certificate is desired explicitly, the finite
period argument in Theorem~\ref{nine-saturation} can be repeated using
the modulus $JDC^2$ instead of $J\sqrt D C^2$.

Conversely, use the inherited canonical ratio witnesses from a source zero.
Set $i=\sqrt D\,i_{\rm old}$ to preserve the old $F,G,H,I$, keeping all
other ratio coordinates fixed.  Choose positive integral roots and put
\[
 n=m_0(V_0+V_1+\sqrt{a_0}+\sqrt{a_1})-W_1^2-Q_0/J.
\]
This is an integer.  The retained bound $W_1^2+Q_0/J<2D$, together with
$V_1>2D$ from $T_0=ADe_0>2D$, proves $n>0$.  One norm factor vanishes.
\end{proof}

Composing this existential equivalence with the named source/refutation
calibration, and then negating, gives the intended ordinary argument in
$\ZFC$.  The new proof does not by itself close the repository's
remaining internal-model calibration or produce a new kernel-checked
object-language consistency biconditional.

\subsection{A smaller-degree majorant of the whole radical}

Temporarily treat $A,C,i,Z$ as independent indeterminates, put $J=4A-5$,
and replace $H$ by $C+FZ$.  Set
\[
 \mathcal Q=DF\bigl((G^2-1)H^2+1\bigr).
\]
Its $A$-degree is 36, with leading coefficient
$(1024C^{19}i^5Z)^2$.  Write $\mathcal Q=\sum q_k A^k$ and define
\begin{align*}
 t_{18}&=1024C^{19}i^5Z,\\
 t_{18-k}&=\frac{q_{36-k}-\sum_{\ell=1}^{k-1}
                  t_{18-\ell}t_{18-k+\ell}}{2t_{18}}
                  \qquad(1\le k\le18).
\end{align*}
This finite calculation takes place in the coefficient field
$\mathbb Q(C,i,Z)$.  It specifies the polynomial part of the formal
square root; it is not a division performed on numerical witnesses.
Exact arithmetic shows that
\[
 \sigma=1024C^5i,\qquad
 \mathcal P=\sigma\sum_{r=0}^{18}t_r A^r
\]
has 333 integer monomials and no negative exponents.  Define
\begin{equation}
 \mathcal L=2^{20}C^{24}i^6Z,\qquad
 \mathcal P_+=\mathcal P+\mathcal L,\qquad
 \mathcal D=\mathcal P_+^2-\sigma^2\mathcal Q.
 \label{nine-radical-polynomial}
\end{equation}
The complete coefficient lists are supplied.  Two exact integer-polynomial
identities give the uniform inequalities:
\begin{itemize}
 \item $\mathcal P(A+3,C+1,i+1,Z+1)$ has 6,650 strictly positive coefficients,
       with constant $15\,915\,642\,730\,652$.
 \item The translated defect $\mathcal D(A+3,C+1,i+1,Z+1)$ has
       36,309 strictly positive coefficients. Its constant is
       $32\,237\,226\,735\,615\,909\,648$.
\end{itemize}
Consequently
\begin{equation}
 \mathcal P_+(A,C,i,Z)>\sigma\sqrt{\mathcal Q}>0
 \quad(A\ge3, C,i,Z\ge1).
 \label{nine-whole-bound}
\end{equation}
This is a universal positivity certificate.  The uncorrected formal
polynomial part would be invalid: its defect at $(3,1,1,1)$ is
$-1\,140\,296\,347\,768\,021\,232$.  The correction in
\eqref{nine-radical-polynomial} is therefore substantive.

For the second fixed equation put
\[
 Z=B_1+(j-1)C,\quad
 V_1=\mathcal P_+(A,C,i,Z),\quad a_1=(1024C^5i)^2DFI,
\]
and use exactly the norm equation~\eqref{nine-cheap-equation}.
The previous proof applies with one additional observation: once the
signed root $r_1$ is integral, $r_1/\sigma$ is a rational square root of
the integer $DFI$, hence is itself an integer.  Thus the scale introduces
no extra square solutions.  In the converse direction $F>D$ and
$\sigma\ge1024$ imply $V_1>2D$, so the witness $n$ remains positive.

The supplied graph evaluates the 333 coefficients by fixed polynomial
operations.  All 24 multivariate Horner orders were counted; the selected
order is $Z,i,C,A$.  This yields 1,074 live binary operations.  It is a
finite evaluated circuit count, not a claim of an optimal arithmetic circuit.

\subsection{Exact degrees and the changed dominant term}

The source has
\[
 r=\deg R=8799,\qquad \nu=\deg N=2324,\qquad W=116,
 \qquad a=r+4\nu+2=18097.
\]
The common register degrees are
\begin{align*}
 \deg\Omega&=5a+2W+3=90720,\\
 \deg V_0&=3a+r+1=63091,\\
 \deg\delta_0&=3a+r+2=63092,\qquad \delta_0=K^2-1,\\
 \deg E_0&=a+4=18101,\qquad\deg a_0=126182.
\end{align*}
For either majorant write $\delta_1=V_1^2-a_1$ and
\[
 B_* =E_0^2-2E_0\Omega(V_0+V_1)+\Omega^2(\delta_0+\delta_1).
\]
The universal exact identity
\begin{equation}
 F_9=B_*^2+4\Omega^2V_0V_1B_*
           +4\Omega^4(\delta_0V_1^2+a_0\delta_1)
 \label{nine-resolved-norm}
\end{equation}
resolves the leading cancellations that a raw syntactic degree count misses.

For the simple majorant,
\[
 \deg V_1=18a+r+24=334569,
 \qquad\deg\delta_1=34a+2r+48=632944.
\]
The latter follows directly from~\eqref{nine-cheap-defect}.  The square
$B_*^2$ is uniquely dominant in~\eqref{nine-resolved-norm}; therefore
\[
 \deg F_9=92r+352\nu+8W+284=1\,628\,768.
\]
Its leading coefficient is positive on the actual source tuple.

For the complete-radical bound the unique greatest weighted monomials,
with weights $(a,1,1,r)$ on $(A,C,i,Z)$, are
\[
 2^{20}A^{18}C^{24}i^6Z\quad\hbox{in }\mathcal P_+,
 \qquad
 2^{41}A^{18}C^{48}i^{12}Z^2\quad\hbox{in }\mathcal D.
\]
Hence $\deg V_1=334575$, $\deg\delta_1=343404$ and
$\deg B_*=524844$.  The four summands in~\eqref{nine-resolved-norm},
splitting its last parenthesis, have degrees
\begin{center}
\begin{tabular}{lr}
$B_*^2$&1,049,688\\
$4\Omega^2V_0V_1B_*$&1,103,950\\
$4\Omega^4\delta_0V_1^2$&1,095,122\\
$4\Omega^4a_0\delta_1$&832,466
\end{tabular}
\end{center}
The second term is now uniquely dominant, with positive leading
coefficient.  Thus the exact degree is
\[
 \boxed{\deg F_9=63r+236\nu+8W+221=1\,103\,950.}
\]
Translation of positive inputs to natural successors preserves these
degrees; the natural $j$ directly supplies the existing $j-1$ coordinate.

\subsection{What the executable evidence proves}

The independent verifier uses a separately implemented sparse integer
polynomial ring and does not import either circuit builder or the rational
discovery code.  It checks the entire radical identity, shifted positive
coefficient lists, all local certificate formulas, the exact final norm,
source-prefix structural equality, topology, liveness, operation counts,
all fixed-leaf hashes and the resolved degree comparisons.  It rejects
both the uncorrected radical truncation and a corrupted final norm.
The complete coefficient generators and parent proofs are retained as
explicit dependencies; no finite check is presented as a consistency test.

The main mathematical change is the saturation theorem, whose soundness
and finite-period completeness were proved above.  The whole-radical
certificate is a second new uniform lemma.  The remaining native/compiler
and syntactic-refutation interfaces are inherited, and their unresolved
formal verification boundary remains visible.

% Integration-ready mathematical appendix. No standalone preamble.
\section{A different source: fixed coprime formula labels}
\label{sec:fermat-flow}

This section supplies a new source-level alternative to the affine-factor
and coefficient-budget sources inspected at commit
\texttt{c1a102936aca437e9999530a0111e6ddab56dc6d}. It does not supply
a smaller completed Diophantine equation. The change removes a freely
chosen common multiple or a coefficient-height budget from the source;
the cost of computing and retrieving its labels remains to be compiled.
The closure argument itself is the single-conclusion Horn-flow lemma
already used by package~68. The different ingredient here is a fixed,
parameter-free sequence of pairwise coprime labels.

\subsection{Coprime labels do not need to be primes}

For $j\in\mathbb N$ put
\[
 F_j=2^{2^j}+1.
\]
The elementary identity
\[
 \prod_{i<j}F_i=F_j-2
\]
follows by induction, including the empty product at $j=0$.
Consequently $F_j>1$ and the $F_j$ are pairwise coprime: a common divisor
of $F_i,F_j$, with $i<j$, divides $F_j$ and $F_j-2$, and is odd.
Primality is unnecessary; $F_5=641\cdot6700417$ is allowed.

For any finitely supported vectors $a,b\in\mathbb N^{(\mathbb N)}$,
\begin{equation}\label{eq:fermat-injective}
 \prod_jF_j^{a_j}=\prod_jF_j^{b_j}
 \quad\Longleftrightarrow\quad a=b.
\end{equation}
To prove the nontrivial implication, cancel the minimum of the two
exponents at every index. If a positive exponent remains on one side,
its factor $F_j>1$ is coprime to the whole other side, a contradiction.
This proof requires no factorization algorithm or primality predicate.

\subsection{An exact unordered proof certificate}

Let $\mathcal R$ be any family of finitary Horn rules with natural-number
labels, one conclusion $c_r$ and a finite premise multiset $P_r$.
Zero-premise rules are axioms. Let $D$ be the least rule-closed set.
Then for each target $t$ the following conditions are equivalent:
\begin{enumerate}
 \item $t\in D$.
 \item Some finite list of valid rule instances satisfies
 \begin{equation}\label{eq:fermat-product}
  \prod_{i<N}F_{c_i}
    =F_t\prod_{i<N}\prod_{p\in P_i}F_p.
 \end{equation}
 \item Some finite list of valid rule instances and natural numbers
 $S_0,\ldots,S_N$ satisfy
 \begin{equation}\label{eq:fermat-accumulator}
  S_0=F_t,\qquad S_N=1,\qquad
  S_{i+1}F_{c_i}=S_i\prod_{p\in P_i}F_p.
 \end{equation}
\end{enumerate}

For completeness, expand a finite proof tree root first. Let $S_i$ be
the product of the weights of the current open obligations. Expanding a
node replaces its conclusion by its actual premises, giving
\eqref{eq:fermat-accumulator}; the empty frontier has product one.

For soundness of the product certificate,
\eqref{eq:fermat-injective} gives the coordinate balance
\[
 \sum_{i<N}\left(e_{c_i}-\sum_{p\in P_i}e_p\right)=e_t.
\]
If $t\notin D$, sum the coordinates outside $D$. Each rule contributes
at most zero: if its conclusion is outside $D$, at least one premise is
outside $D$, while a conclusion inside $D$ contributes zero minus a
nonnegative number. The right side is one, a contradiction. Thus cycles
cannot create a false proof. A disconnected cycle can be present in a
valid certificate, but it is harmless. Saturating the finitely many
selected rules gives a finite proof DAG for the target; this can be
unfolded to a proof tree.

The accumulator conditions imply the product certificate by multiplying
and cancelling. All $S_i$ are positive: this follows inductively from
$S_0>0$ and positivity of every factor in
\eqref{eq:fermat-accumulator}. No independent positivity guard is needed
when the accumulators range over naturals. We do not claim that an
arbitrary ordering of an arbitrary balanced multiset gives an integral
accumulator path. Existence is enough: the recovered proof tree provides
another valid list and path if necessary.

Exactly one conclusion and nonnegative rule multiplicities are essential.
The rule $a\longrightarrow\{a,t\}$ has balance $e_t$ without any axiom.
Also, subtracting the vector for $a,a\longrightarrow b$ from that for
$a\longrightarrow b$ gives $e_a$, although neither rule derives anything
from the empty set. Those are excluded by the stated hypotheses.

For the repository's target-one calculus, use the actual arity of each
rule, so that truth code zero is a genuine zero-premise axiom. Then
$F_0=3$ and the initial target weight is the fixed integer $F_1=5$.
One must not simultaneously use $F_0=3$ and treat padded truth premises
as free factors. A different, also sound convention sets the zero-label
weight to one and retains Fermat weights only for positive labels, but
requires a separate zero-label branch. The convention here avoids it.

\subsection{A precise four-parameter bounded source}

The double exponent in $F_j$ is not an ordinary polynomial. To make its
remaining compilation cost concrete, define
\[
 \beta(x,y,j)=x\bmod(1+(j+1)y).
\]
Use four outer natural parameters $N,x_S,x_Q,y$, and put
$S_i=\beta(x_S,y,i)$, $Q_i=\beta(x_Q,y,i)$.
Impose the endpoints
\[
 S_0=5,\qquad S_N=1,\qquad Q_0=2.
\]
For every $i<N$, require simultaneously
\[
 Q_{i+1}=Q_i^2
\]
and one actual rule instance $P_i\longrightarrow c_i$, all of whose
labels are strictly less than $N$, with
\[
 S_{i+1}\bigl(\beta(x_Q,y,c_i)+1\bigr)
   =S_i\prod_{p\in P_i}\bigl(\beta(x_Q,y,p)+1\bigr).
\]
Allow identity rules $c\longrightarrow c$ as padding; their addition
does not change derivability. These conditions use one bounded universal
block, with existential local witnesses. Induction gives
$Q_j=2^{2^j}$ for $0\le j\le N$, so every referenced label is authenticated.
The preceding theorem proves soundness.

For completeness, choose a finite proof tree and then a horizon $N$
larger than all labels and at least its number of nodes. Pad the rule list
by identities, keeping the terminal accumulator equal to one. Form the
finite $S$ and $Q$ sequences first. Choose
\[
 y=N!\left(1+\max_{0\le i\le N}\{S_i,Q_i\}\right).
\]
The moduli $1+(i+1)y$ are pairwise coprime and exceed the desired entries.
The finite Chinese remainder theorem supplies $x_S,x_Q$ with the same
scale $y$. This factorial is used only to prove existence of CRT witnesses;
it is not a function in the proposed final polynomial, and no source
condition asserts that an unknown is a factorial. The choice order avoids
any self-referential size bound.

\subsection{What must still be charged}

The preceding four-parameter source is a finite first-order arithmetic
formula, not an ordinary fixed-dimensional polynomial. Its costs include
the supplied Horn-rule recognizer, the two adjacent sequence entries,
up to three indexed label lookups per binary rule, label-range witnesses,
and the bounded-universal compiler. For example a lookup
$u=\beta(x_Q,y,p)$ has the completely explicit quadratic graph
\[
 d=(p+1)y,\qquad x_Q=(d+1)k+u,\qquad u+s=d
\]
over naturals. These are three equations using four local witnesses
$d,k,u,s$, one of which is the returned value $u$. A strict label bound
is $p+b+1=N$ with one further natural witness $b$.
Shared denominators and fixed zero-premise branches may save registers,
but no unmeasured savings are assumed.

Alternatively, one can define a label by two separately compiled
exponentiation graphs. That removes the power sequence and introduces the
corresponding graph witnesses. Both routes remove the source's common
multiple or coefficient budget; neither has yet been shown to reduce the
complete polynomial's variable count, degree or arithmetic cost.
The finite script \texttt{alternative/verify.py} checks all rule-multiset cases
listed in \texttt{alternative/VERIFICATION.json}, constructs proof trees and
integer traces, verifies a small shared-CRT example, and includes the
failure modes above. It is an adversarial regression, not a substitute
for these universal proofs or the inherited calculus-to-ZFC calibration.


\section{Reproducibility, source binding, and integration}

\subsection{What the archive contains}

The main deliverables are this article and its editable LaTeX source.
The \file{integer/} directory contains four explicit equations, paid
circuits, six fixed coefficient payloads and their recipes, source copies,
the inverse maps, and independent polynomial checks. The \file{nine/}
directory contains both positive constructions, natural translations,
the complete fixed-leaf binding, source/layout/kernel data, and the
radical-majorant certificates. The \file{alternative/} directory contains
the separate finite proof-flow construction and its diagnostic examples.
\file{provenance/BASELINE.json} pins the reviewed commit and all five
incoming ZIPs by SHA-256 and Git blob identifier. Component notices retain
the inherited licenses and attribution.

Use Python 3.10 or newer for the standard-library checks:
\begin{lstlisting}
python3 verify.py
\end{lstlisting}
The verification receipt distinguishes exact algebra from finite
diagnostics and the inherited unbounded proof dependencies. The optional
discovery calculations are separate from checking the released polynomial
and its certificates. The full coefficient bytes are included, so validation
does not require a network connection or a live GitHub branch.

To rebuild this PDF with a standard LaTeX installation:
\begin{lstlisting}
make pdf
\end{lstlisting}
To rebuild the release ZIP and its content manifest:
\begin{lstlisting}
python3 make_archive.py
\end{lstlisting}

\subsection{What should be incorporated into ProveIt}

The most direct new candidate is the 18-integer degree-88 equation: its
proof interface is a short polynomial automorphism and monic elimination.
It improves a recently delivered low-degree branch without changing its
coefficient tuple. The three lower-degree companions should be retained
with it rather than choosing a single point without a cost criterion.

For the nine-positive branch, the saturated short-majorant circuit improves
both cost and degree over the current 440-operation candidate. The
complete-radical form should be retained as a second candidate when total
degree is the priority. Its large positivity certificate is a finite
proof object that can be checked independently; it should not be confused
with a short hand-expanded formula or a proof of the native compiler.

The direct proof-flow construction belongs with source-level alternatives.
Its removal of a global proof-order certificate is established, but the
lookup and aggregation costs must be compiled and counted before it can
replace a finished consistency equation. These are concrete next
obligations, not missing parameters hidden inside the reported metrics.

\subsection{Limits of the result}

No global simplicity minimum is proved. The fewest possible number of
variables for a polynomial specific to this consistency statement is not
determined. We do not infer historical novelty from the absence of a formula
in the inspected snapshot. Large fixed numerals, high total degree, short
arithmetic circuits, and low witness counts remain different objectives.

No zero search or consistency assertion is part of verification. The
equations are designed so that any zero provides a finite refutation and
any finite refutation supplies a zero through the stated proof chain.
Their mathematical significance is the biconditional, including the
reverse construction, not an empirical failure to find a solution.

\begingroup
\small
\begin{thebibliography}{99}
\raggedright
\bibitem{repo}
Vladimir Reshetnikov and contributors.
\emph{ProveIt, DiophantineConsistency research directory}, snapshot
\texttt{c1a102936aca437e9999530a0111e6ddab56dc6d}.
\url{https://github.com/VladimirReshetnikov/ProveIt/tree/c1a102936aca437e9999530a0111e6ddab56dc6d/Computability/DiophantineConsistency}.

\bibitem{incoming}
ProveIt incoming research reports, same snapshot.
\url{https://github.com/VladimirReshetnikov/ProveIt/tree/c1a102936aca437e9999530a0111e6ddab56dc6d/docs/incoming}.
Individual archive hashes appear in \file{provenance/BASELINE.json}.

\bibitem{source73}
ProveIt source 73.
\emph{Reanchored sign constraints and adaptive Pell carriers},
section ``Exact local equivalence,'' 2026.
\url{https://github.com/VladimirReshetnikov/ProveIt/tree/c1a102936aca437e9999530a0111e6ddab56dc6d/Computability/DiophantineConsistency/con-zfc-diophantine-forms/73-guard-reanchoring}.

\bibitem{source66}
ProveIt source 66 and compact54-sharp candidate.
\emph{A refined free-carrier nine-positive equation}; retained source,
packing, ratio, and native coefficient proofs. Snapshot as in \cite{repo}.

\bibitem{joint}
\emph{Joint guards and Pell carriers for Diophantine consistency}.
Incoming \file{ZFC_Diophantine_Joint_Guards_2026-10-24.zip};
\file{joint/PROOF.md}, \file{joint/COMPACT_BASE24.json}, and
the 19-variable degree-88 projection. Pinned in \cite{incoming}.

\bibitem{round11}
\emph{Structural reductions for Diophantine consistency equations}.
Incoming \file{ConZFC_Structural_Reductions_Round11_2026-10-24.zip},
Lucas norm-saturation components. Pinned in \cite{incoming}.

\bibitem{radicalprior}
\emph{Radical majorants, guard compression, and Pell descent}.
Incoming \file{con-zfc-radical-guard-descent-2026-10-23.zip},
nine-positive degree-1,158,228 component. Pinned in \cite{incoming}.

\bibitem{sun2026}
Zhi-Wei Sun.
\emph{Ten unknowns for Hilbert's tenth problem over the integers}.
arXiv:2609.01594v1, 2026.
\url{https://arxiv.org/abs/2609.01594}.

\bibitem{bayer2025}
Jonas Bayer, Marco David, Malte Ha{\ss}ler, Yuri Matiyasevich,
and Dierk Schleicher.
\emph{Diophantine equations over $\Z$: universal bounds and parallel
formalization}. arXiv:2506.20909v2, 2025.
\url{https://arxiv.org/abs/2506.20909}.

\bibitem{jones1982}
James P. Jones.
Universal Diophantine equation.
\emph{Journal of Symbolic Logic} 47(3), 549--571, 1982.
\url{https://doi.org/10.2307/2273588}.
\end{thebibliography}
\endgroup
\end{document}
